\chapter{Contrats à terme datés, options et ETF}\label{ch:m3:dated-futures-options-and-etfs}

Le 10 janvier 2024, la Securities and Exchange Commission (SEC) américaine laisse coter des \omterm{def:m3:dated-futures-options-and-etfs:etp}{produits négociés en bourse} (ETP) adossés au bitcoin au comptant, après qu'un tribunal
a jugé que celle-ci n'avait pas suffisamment motivé son refus à l'égard de l'un d'eux. Le
lendemain, leurs parts commencent à se négocier sur les bourses américaines. D'ici septembre 2026, la plus importante en détenait environ
67~milliards de dollars de bitcoin pour ses porteurs, et un trader pouvait détenir du bitcoin sur un compte de courtage, le prêter, le financer
et vendre contre lui un contrat à terme réglementé. Les produits dérivés du marché crypto se répartissent désormais en deux familles~: les
\omterm{def:m3:dated-futures-options-and-etfs:dated}{contrats margés en crypto} des plateformes natives de la crypto, avec leurs propres conventions, et les contrats margés en dollars
et les fonds des marchés réglementés du livre~1. Ce chapitre traite des contrats à terme datés sur
les deux, du basis trade entre elles, des options de la plateforme d'options crypto dominante, avec leur prime
payée en coin, la forme de la surface de volatilité crypto, et les \omterm{def:m3:dated-futures-options-and-etfs:etp}{produits négociés en bourse}, ainsi que la manière dont leurs
parts sont créées.

\section{Contrats à terme datés sur les plateformes crypto et sur les bourses réglementées}

\begin{definition}[Contrat à échéance, contrat margé en crypto]\label{def:m3:dated-futures-options-and-etfs:dated}
Un \emph{contrat à échéance}\index{contrat à échéance} sur un actif crypto est un contrat à terme à échéance fixe,
réglé en espèces, à l'échéance, sur un indice ou un taux de référence, à la différence d'un \omterm{def:m3:perpetual-futures:perp}{contrat perpétuel}.
Un \emph{contrat margé en crypto}\index{contrat margé en crypto} est un contrat à terme ou d'options dont la marge,
le P\&L et le règlement sont libellés dans le coin sous-jacent plutôt qu'en dollars ou en \omterm{def:m3:blockchains-for-traders:stable}{stablecoin}.
\end{definition}

Les plateformes natives de la crypto cotent les \omterm{def:m3:dated-futures-options-and-etfs:dated}{contrats à échéance} comme des \omterm{def:m3:perpetual-futures:types}{contrats inverses} (\cref{ch:m3:perpetual-futures})~: un notionnel
en dollars, marge et règlement en coin, et des échéances hebdomadaires, mensuelles et trimestrielles, toutes à la même heure.
Les bourses réglementées cotent des contrats réglés en dollars sur le modèle de tout autre contrat à terme sur indice réglé en espèces (livre~1,
chapitre~18), réglés sur un indice de référence calculé à partir de transactions au comptant.

\begin{definition}[Taux de référence crypto]\label{def:m3:dated-futures-options-and-etfs:rate}
Un \emph{taux de référence crypto}\index{taux de référence crypto} est un prix de référence d'un actif crypto calculé par un
administrateur à partir de transactions sur un ensemble de plateformes de négociation au comptant constitutives, sur une fenêtre de temps déterminée, avec des règles (partitions,
médianes, critères sur les plateformes) qui la rendent coûteuse à déplacer avec quelques transactions~; elle sert à régler les contrats à terme et à valoriser les
fonds.
\end{definition}

\begin{dated}{2026-09}{Un taux de référence et ses usages}\label{dat:m3:dated-futures-options-and-etfs:brr}
Le CME CF Bitcoin Reference Rate est un indice de référence calculé une fois par jour, qui agrège les transactions des marchés bitcoin--dollar sur
les bourses constitutives, et il sert d'indice de règlement aux contrats à terme sur bitcoin cotés par CME Group, qui les a autocertifiés le
1\up{er} décembre 2017, pour une première date de négociation fixée au 18 décembre. Sa variante new-yorkaise,
utilisée pour valoriser le plus grand fonds bitcoin au comptant, retient chaque transaction au comptant sur ses plateformes constitutives (citées dans le
rapport annuel 2025 du fonds~: Bitstamp, Coinbase, itBit, Kraken, Gemini, LMAX Digital, Crypto.com et Bullish)
entre 15~h et 16~h, heure de New York, les répartit en douze partitions de cinq minutes, retient pour chaque partition le
prix médian pondéré par les volumes, puis moyenne les douze médianes à pondérations égales.
\end{dated}

La construction est choisie contre la manipulation~: un trader qui passe une grande transaction isolée à un prix hors marché ne
déplace guère la médiane d'une seule partition, et doit déplacer la médiane de toutes les plateformes constitutives pendant une heure entière pour
faire bouger fortement le taux. La construction de ce chapitre la met en œuvre.

Les contrats à terme des deux familles de plateformes ne diffèrent pas seulement par la marge. Le P\&L en dollars d'un \omterm{def:m3:dated-futures-options-and-etfs:dated}{contrat margé en crypto} est le
P\&L en coin multiplié par le cours du coin, de sorte qu'un trader en dollars acheteur du \omterm{def:m3:dated-futures-options-and-etfs:dated}{contrat margé en crypto} est doublement acheteur, comme
avec les \omterm{def:m3:perpetual-futures:perp}{contrats perpétuels} inverses. Le contrat à terme réglementé est margé en dollars par l'intermédiaire d'un membre compensateur, dans la
cascade du livre~1, chapitre~5, et non la cascade au niveau des plateformes du \cref{ch:m3:margin-liquidation-and-loss-allocation}.

\section{Le basis trade sur contrats à terme réglementés}

Un \omterm{def:m3:dated-futures-options-and-etfs:dated}{contrat à échéance} se négocie avec une base par rapport au comptant. Avec $F$ le contrat à terme, $S$ l'indice au comptant et $\tau$ le nombre de jours avant l'échéance, la
base annualisée est $(F/S - 1)\times 365/\tau$ (ou $\ln(F/S)\times 365/\tau$ en capitalisation continue). Le \omterm{def:m3:forward-curves-storage-and-convenience-yield:carry}{cash-and-carry}
(livre~1, chapitre~21) consiste à acheter au comptant, à vendre le contrat à terme et à encaisser la base s'il est conservé jusqu'à l'échéance.

\begin{dated}{2026-09}{Une courbe de contrats à terme crypto}\label{dat:m3:dated-futures-options-and-etfs:curve}
Le 24 septembre 2026, l'indice bitcoin de Deribit valait 83\,499.12 et ses contrats à terme trimestriels étaient valorisés à 84\,533.98
(décembre 2026), 85\,668.78 (mars 2027), 86\,833.65 (juin 2027) et 87\,948.30 (septembre 2027)~: base annualisée de
4.9~\%, 5.2~\%, 5.3~\% et 5.3~\%. Ses contrats à terme et options expirent à 08:00 UTC.
\end{dated}

\begin{proposition}[Le carry d'un basis trade fonds--contrat à terme]\label{prop:m3:dated-futures-options-and-etfs:carry}
On achète au comptant un \omterm{def:m3:dated-futures-options-and-etfs:etp}{produit négocié en bourse} (ETP) financé au taux $r_{\mathrm{fin}}$, on paie ses frais annuels $\phi_{\mathrm{ETF}}$,
on vend un contrat à terme réglementé à une base annualisée $b$, avec des coûts de contrats à terme $c$ par an (frais et roll). Si le fonds
suit le comptant et que la position est conservée jusqu'à l'échéance du contrat à terme, le carry par unité de notionnel est
$b - r_{\mathrm{fin}} - \phi_{\mathrm{ETF}} - c$ par an, et son rendement sur le capital immobilisé (la marge du contrat à terme $m$ plus
la décote $h$ sur le fonds financé) est ce carry divisé par $m + h$.
\end{proposition}

\begin{proof}
À l'échéance, le contrat à terme converge vers le taux de référence, et les parts du fonds, par création et rachat, vers leur valeur
liquidative, qui suit le même prix au comptant~; les deux jambes de prix s'annulent et ce qui subsiste est la base encaissée, moins les coûts de
détention du fonds. Le capital correspond à ce que retiennent le financeur et la chambre de compensation.
\end{proof}

La stratégie ne vaut que ce que vaut la base nette de financement. \Cref{tab:m3:dated-futures-options-and-etfs:carry} compare une
base élevée de 10~\% par an, niveau moyen du carry crypto mesuré par les chercheurs de la BRI
(\cref{ch:m3:perpetual-futures}), avec les 4.9~\% de l'instantané ci-dessus, pour un \omterm{def:m3:perpetual-futures:perp}{taux de financement} illustratif de 4.5~\%. À une
base élevée, la stratégie rapporte environ 5~\% par an sur le notionnel et 14~\% sur le capital~; avec la base de l'instantané, elle perd de l'argent. Les
risques de la stratégie sont ceux d'un carry trade à effet de levier~: la base peut s'élargir contre le contrat vendu avant l'échéance, ce qui exige
une marge supplémentaire~; le financement peut être retiré~; et le fonds peut se traiter à une décote par rapport à sa valeur liquidative lorsque les créations et
rachats sont entravés.

\begin{table}[ht]
\centering
\small
\begin{tabular}{@{}lrr@{}}
\toprule
par an, en~\% du notionnel & base élevée & base de l'instantané \\
\midrule
Base encaissée sur le contrat vendu & 10.00 & 4.92 \\
Financement du fonds & $-4.50$ & $-4.50$ \\
Frais du fonds & $-0.25$ & $-0.25$ \\
Coûts des contrats à terme & $-0.30$ & $-0.30$ \\
\midrule
Carry & 4.95 & $-0.13$ \\
Rendement du capital (marge de 25~\% + décote de 10~\%) & 14.14 & $-0.38$ \\
\bottomrule
\end{tabular}
\caption{Le carry d'un fonds au comptant en position longue et d'un contrat à terme réglementé en position courte (\cref{prop:m3:dated-futures-options-and-etfs:carry}). La
base élevée est la moyenne de la BRI pour le carry crypto~; l'instantané correspond au contrat à terme de décembre 2026 de
l'\cref{dat:m3:dated-futures-options-and-etfs:curve}~; le financement, les coûts, la marge et la décote sont illustratifs~; les frais sont ceux
du plus grand fonds.}\label{tab:m3:dated-futures-options-and-etfs:carry}
\end{table}

\section{Les conventions de la plateforme d'options dominante~: les options inverses}

Sur les plateformes natives de la crypto, les options suivent les conventions de Deribit, décrites ci-dessous~; les bourses réglementées cotent des
options réglées en dollars sur contrats à terme, et les options sur les fonds au comptant se négocient sur les bourses de valeurs, la première ayant été approuvée par la SEC
le 20 septembre 2024 pour l'iShares Bitcoin Trust.

\begin{definition}[Option inverse]\label{def:m3:dated-futures-options-and-etfs:inverse}
Une \emph{option inverse}\index{option inverse} est une option européenne sur le prix en dollars d'une crypto-monnaie, dont la prime est
cotée et payée en crypto-monnaie et dont le gain est réglé en crypto-monnaie~: un call verse $(F_T - K)^+/F_T$ unités de la crypto-monnaie à l'échéance, où
$F_T$ est le prix de règlement et $K$ le prix d'exercice en dollars.
\end{definition}

\begin{dated}{2026-09}{Conventions d'options sur la plus grande plateforme d'options crypto}\label{dat:m3:dated-futures-options-and-etfs:deribit}
Les données d'instruments de Deribit au 24 septembre 2026 décrivent ses options sur bitcoin comme européennes, d'un bitcoin par contrat,
cotées et réglées en bitcoin, avec des prix d'exercice en dollars et des échéances à 08:00 UTC, quotidiennes, hebdomadaires, mensuelles et trimestrielles. Ce
jour-là, les options de décembre 2026, dont le contrat à terme sous-jacent valait 84\,538.67, étaient valorisées à des volatilités implicites de
47.1~\% (prix d'exercice 60\,000), 40.8~\% (70\,000), 38.2~\% (80\,000), 37.5~\% (90\,000), 38.1~\% (100\,000), 39.6~\% (110\,000) et
42.2~\% (120\,000). Coinbase, qui a acquis Deribit le 14 août 2025 pour environ 4.3~milliards de dollars, la décrit dans son
rapport annuel 2025 comme le premier acteur mondial des options crypto en volume échangé et en positions ouvertes.
\end{dated}

\begin{proposition}[Prix et deltas des options inverses]\label{prop:m3:dated-futures-options-and-etfs:inverse}
Soit $V$ le prix de Black en dollars d'une option sur le contrat à terme $F$ (livre~1, chapitre~25), sans actualisation.
La prime de l'\omterm{def:m3:dated-futures-options-and-etfs:inverse}{option inverse} en unités de crypto-monnaie vaut $V/F$~; la parité call-put en unités de crypto-monnaie s'écrit $C/F - P/F = 1 - K/F$. Le delta de l'option en
unités de crypto-monnaie, pour un porteur qui a payé la prime en crypto-monnaie, est le delta en dollars $\partial V/\partial F$ diminué de la prime en crypto-monnaie $V/F$~: le
delta ajusté de la prime. À l'échéance, la valeur en crypto-monnaie d'un call, $(F_T - K)^+/F_T$, reste inférieure à un, quelle que soit la hausse de $F_T$,
alors que celle d'un put, $(K - F_T)^+/F_T$, croît sans borne lorsque $F_T \to 0$.
\end{proposition}

\begin{proof}
La valeur en dollars $V$ vaut $V/F$ unités de crypto-monnaie au cours actuel~; en divisant la parité de Black $C - P = F - K$ par $F$, on obtient la
forme en crypto-monnaie. Un porteur dont le compte est en crypto-monnaie et qui a acheté l'option avec $V/F$ unités dispose d'une exposition en dollars
égale à $\partial V/\partial F$ provenant de l'option et à $-V/F$ provenant des unités dépensées~; sa couverture en crypto-monnaie est la somme des deux. Les bornes
découlent de $(F_T - K)/F_T = 1 - K/F_T < 1$ et de $K/F_T - 1 \to \infty$.
\end{proof}

La convention est celle du delta ajusté de la prime pour les options de change (livre~2, chapitre~19)~: la prime est payée dans
l'actif dont le prix sert de sous-jacent à l'option. Les propres cotations de Deribit confirment la tarification~: la cotation de 0.0492 bitcoin du call de décembre 90\,000 correspond
au prix de Black en dollars, calculé à une volatilité de 37.46~\%, divisé par le contrat à terme.

\begin{omfigure}
\begin{tikzpicture}
\begin{axis}[width=11cm,height=5cm,
  xlabel={prix du bitcoin à l'échéance, USD},ylabel={valeur à l'échéance, BTC},
  tick label style={font=\small},label style={font=\small},
  xmin=30000,xmax=180000,ymin=0,ymax=2.1,grid=major,grid style={gray!25},
  xtick={30000,60000,90000,120000,150000,180000},xticklabels={30\,000,60\,000,90\,000,120\,000,150\,000,180\,000},scaled x ticks=false,
  legend columns=2,legend style={font=\footnotesize,at={(0.5,-0.3)},anchor=north,draw=none,fill=none,/tikz/every even column/.append style={column sep=8pt}},
  table/col sep=comma]
\addplot[omThm,very thick] table[x=price,y=call_coin]{figdata/markets-3/19-dated-futures-options-and-etfs/expiry.csv};\addlegendentry{call inverse, prix d'exercice 90\,000}
\addplot[omDef,very thick,dashed] table[x=price,y=put_coin]{figdata/markets-3/19-dated-futures-options-and-etfs/expiry.csv};\addlegendentry{put inverse, prix d'exercice 90\,000}
\end{axis}
\end{tikzpicture}

\omcaption{Valeur en bitcoin à l'échéance d'\omterm{def:m3:dated-futures-options-and-etfs:inverse}{options inverses} de prix d'exercice 90\,000 (\cref{prop:m3:dated-futures-options-and-etfs:inverse})~:
la valeur en bitcoin du call s'approche d'un bitcoin sans jamais l'atteindre~; celle du put croît sans borne lorsque le cours baisse. En dollars,
les deux sont les gains ordinaires. Données~: le tutoriel du chapitre.}\label{fig:m3:dated-futures-options-and-etfs:expiry}
\end{omfigure}

\begin{omfigure}
\begin{tikzpicture}
\begin{axis}[width=11cm,height=5cm,
  xlabel={prix d'exercice, milliers de dollars},ylabel={delta du call},
  tick label style={font=\small},label style={font=\small},
  xmin=60,xmax=120,ymin=0,ymax=1,grid=major,grid style={gray!25},
  legend columns=2,legend style={font=\footnotesize,at={(0.5,-0.3)},anchor=north,draw=none,fill=none,/tikz/every even column/.append style={column sep=8pt}},
  table/col sep=comma]
\addplot[omThm,very thick,mark=*,mark size=1.5pt] table[x=strike,y=forward]{figdata/markets-3/19-dated-futures-options-and-etfs/deltas.csv};\addlegendentry{delta en dollars (à terme)}
\addplot[omDef,very thick,dashed,mark=square*,mark size=1.5pt] table[x=strike,y=adjusted]{figdata/markets-3/19-dated-futures-options-and-etfs/deltas.csv};\addlegendentry{delta ajusté de la prime}
\end{axis}
\end{tikzpicture}

\omcaption{Deltas des calls inverses de décembre 2026 selon le prix d'exercice, aux volatilités marquées par la plateforme de
l'\cref{dat:m3:dated-futures-options-and-etfs:deribit}. Dans la monnaie profonde, le delta ajusté de la prime baisse~: la prime, payée
en bitcoin, représente elle-même une forte exposition au bitcoin, et le delta ajusté du call 60\,000 (0.645) est inférieur à celui du call 70\,000 (0.658).
Données~: le tutoriel du chapitre.}\label{fig:m3:dated-futures-options-and-etfs:deltas}
\end{omfigure}

\section{La surface de volatilité des crypto-monnaies}

L'instantané de l'\cref{dat:m3:dated-futures-options-and-etfs:deribit} présente la forme que prend habituellement la surface crypto~: des niveaux élevés
(la volatilité implicite à la monnaie du bitcoin, voisine de 37~\% lors d'une journée calme), un sourire dont les deux ailes sont cotées à l'achat, et un skew dont le signe peut
changer avec l'humeur du marché. Dans l'instantané, l'aile située 30~\% sous le contrat à terme porte environ 10 points de volatilité de plus que
l'aile située 40~\% au-dessus~: un skew de puts. Lors des phases de hausse spéculative, l'aile des calls peut devenir la plus chère, les acheteurs payant pour la hausse~; la
surface des indices actions du livre~1, chapitre~25, présente presque toujours un skew de puts. Le marché se négocie en continu, week-ends
compris, si bien qu'il n'existe ni clôture nocturne ni clôture de week-end, et le temps jusqu'à l'échéance se mesure naturellement en temps calendaire.
Les outils de la surface (paramétrisations du sourire, interpolation sans arbitrage du livre~5) s'appliquent sans modification dès que les primes sont
converties en dollars grâce à la \cref{prop:m3:dated-futures-options-and-etfs:inverse}.

\section{ETF au comptant~: création et rachat}

\begin{definition}[Produit négocié en bourse, grantor trust]\label{def:m3:dated-futures-options-and-etfs:etp}
Un \emph{produit négocié en bourse}\index{produit négocié en bourse} (ETP) est un titre coté en bourse dont les parts suivent un actif ou un indice
et qui sont créées et rachetées par des participants autorisés, qu'il soit ou non enregistré comme fonds de placement~; aux États-Unis,
les produits au comptant sur bitcoin sont des trusts qui détiennent du bitcoin et ne sont pas des sociétés d'investissement enregistrées. Un \emph{grantor
trust}\index{grantor trust} est un trust dont les propriétaires sont traités, pour l'impôt fédéral américain sur le revenu, comme détenant directement ses actifs, de sorte que le
trust lui-même n'est pas imposé.
\end{definition}

\begin{definition}[Création en espèces]\label{def:m3:dated-futures-options-and-etfs:cash}
La \emph{création en espèces}\index{création en espèces} est la création de parts d'un \omterm{def:m3:dated-futures-options-and-etfs:etp}{produit négocié en bourse} contre des espèces remises par un
participant autorisé, le fonds achetant lui-même l'actif sous-jacent~; le rachat en espèces en est l'inverse. Elles s'opposent à la
création et au rachat en nature, dans lesquels le participant livre ou reçoit l'actif.
\end{definition}

Le mécanisme est celui du fonds indiciel coté du livre~1, chapitre~14~: les participants autorisés créent des paniers de parts lorsque
les parts se négocient avec une prime sur la valeur liquidative et les rachètent avec une décote. Pour un trust bitcoin, l'actif est une crypto-monnaie détenue auprès
d'un dépositaire, et la question de savoir qui la manipule détermine l'efficacité de l'arbitrage. Avec la \omterm{def:m3:dated-futures-options-and-etfs:cash}{création en espèces}, le participant ne
détient jamais de bitcoin~: le trust en achète, au taux de référence ou par l'intermédiaire d'un agent d'exécution, et l'arbitrage du participant porte sur l'écart entre le cours de la
part et un montant en espèces fixé ultérieurement, ce qui est plus difficile à couvrir. Avec la création en nature, le participant livre du bitcoin, peut
se le procurer là où il est le moins cher, et se couvre exactement.

\begin{dated}{2026-09}{Le plus grand fonds bitcoin au comptant}\label{dat:m3:dated-futures-options-and-etfs:ibit}
Le rapport annuel 2025 de l'iShares Bitcoin Trust indique que son enregistrement est devenu effectif le 10 janvier 2024, que ses
parts ont été admises à la cote du Nasdaq le 11 janvier 2024, que le fonds crée et rachète par paniers de 40\,000 parts via des participants autorisés,
que ses bitcoins sont conservés par Coinbase Custody Trust Company, que le fonds valorise ses bitcoins selon la variante new-yorkaise du CME
CF Bitcoin Reference Rate, et que, le 29 juillet 2025, la SEC a émis des ordonnances autorisant pour le trust la création et le rachat en nature~;
ses frais de sponsor s'élevaient à 0.25~\% par an, et le rapport suppose qu'il est un \omterm{def:m3:dated-futures-options-and-etfs:etp}{grantor trust} au regard de la fiscalité américaine. Le site du fonds indiquait un actif net de
67.08~milliards de dollars au 23 septembre 2026 et une prime sur la valeur liquidative de $-0.21\%$ le 22 septembre.
\end{dated}

\begin{omfigure}
\centering
\begin{tikzpicture}[font=\small,
  p/.style={draw,thick,rounded corners=2pt,align=center,minimum height=1cm,text width=2.6cm},
  arr/.style={-{Stealth[length=2.2mm]},thick}]
\node[p,draw=omDef,fill=omDef!8] (ap) at (0,0) {participant\\autorisé};
\node[p,draw=omThm,fill=omThm!8,text width=3.1cm] (tr) at (5.6,0) {trust\\(valeur liquidative\\au taux de référence)};
\node[p,draw=omProp,fill=omProp!8] (cu) at (10.9,0) {dépositaire\\(bitcoin)};
\node[p,draw=omExo,fill=omExo!8] (mk) at (5.6,-2.6) {marché au comptant\\(agent d'exécution)};
\draw[arr] ([yshift=8pt]ap.east) -- node[above,font=\footnotesize] {espèces ou bitcoin} ([yshift=8pt]tr.west);
\draw[arr] ([yshift=-8pt]tr.west) -- node[below,font=\footnotesize] {panier de parts} ([yshift=-8pt]ap.east);
\draw[arr] (tr) -- node[above,font=\footnotesize] {bitcoin} (cu);
\draw[arr,dashed] (tr) -- node[right,font=\footnotesize,align=left] {création en espèces~:\\le trust achète} (mk);
\end{tikzpicture}

\omcaption{Création de parts dans un trust bitcoin au comptant. En nature, le participant autorisé livre du bitcoin, qui est transféré au dépositaire~;
en espèces, le participant livre des dollars et le trust achète le bitcoin par l'intermédiaire d'un agent d'exécution. Le rachat inverse les flèches. Schéma de principe.}\label{fig:m3:dated-futures-options-and-etfs:creation}
\end{omfigure}

Les fonds ont modifié l'infrastructure du marché. Le bitcoin détenu pour leurs porteurs de parts est confié à quelques dépositaires~; les créations et les rachats déplacent
des bitcoins par blocs importants à la clôture new-yorkaise~; et le basis trade de la \cref{prop:m3:dated-futures-options-and-etfs:carry} a relié le
contrat à terme réglementé au fonds, de sorte qu'une base élevée attire des créations et une base qui s'effondre attire des rachats. Un trust coté qui ne peut pas racheter ses parts
contre ses actifs laisse le cours des parts s'écarter de leur valeur~: de mai 2015 à décembre 2022, le trust bitcoin de Grayscale s'est négocié entre une
prime de 142~\% et une décote de 49~\% par rapport à la valeur de son bitcoin par part, et a terminé 2022 avec une décote de 45~\%. La décision de justice qui a précédé les
autorisations de 2024 concernait sa demande de conversion en \omterm{def:m3:dated-futures-options-and-etfs:etp}{produit négocié en bourse} (ETP).

\section{Tutoriel~: primes, deltas et base}\label{tut:m3:dated-futures-options-and-etfs:convert}

\begin{tutorial}
\textbf{Objectif.} Convertir les primes des \omterm{def:m3:dated-futures-options-and-etfs:inverse}{options inverses} entre bitcoin et dollars, calculer le delta qui intègre la devise de la prime, annualiser
la base des \omterm{def:m3:dated-futures-options-and-etfs:dated}{contrats à échéance} et chiffrer le coût du trade fonds--contrat à terme. \textbf{État final~:}
les \cref{fig:m3:dated-futures-options-and-etfs:expiry,fig:m3:dated-futures-options-and-etfs:deltas} et
le \cref{tab:m3:dated-futures-options-and-etfs:carry}.
\begin{tutsteps}
\item \textbf{Primes et deltas.} Le prix de Black sur le contrat à terme, divisé par le contrat à terme.
\omcode{code/firm/cryptoopt/firm_cryptoopt.py}{16}{35}{Prime d'option inverse en bitcoin, ainsi que le delta en dollars et le delta ajusté de la prime.}\label{lst:m3:dated-futures-options-and-etfs:delta}
\item \textbf{Le taux de référence.} Douze partitions, une médiane pondérée par les volumes dans chacune, et leur moyenne.
\omcode{code/firm/cryptoopt/firm_cryptoopt.py}{67}{80}{Un taux de référence dans le style d'un indice bitcoin publié.}\label{lst:m3:dated-futures-options-and-etfs:brr}
\item \textbf{La base et le carry.} \texttt{basis\_table()} et \texttt{scenarios()}~; \texttt{fig\_etf.py} écrit les données du graphique.
\end{tutsteps}
\textbf{Pour aller plus loin.} Recalculer les deltas avec la pente du sourire de volatilité (le «~delta de sourire~»)~; évaluer un put spread en bitcoin et en dollars
et comparer leurs couvertures~; tester combien de transactions d'une minute un attaquant doit exécuter pour déplacer le taux de référence de 1~\%.
\end{tutorial}

\section{Construction~: options inverses et base}\label{bld:m3:dated-futures-options-and-etfs:cryptoopt}

\begin{build}
\bfield{Rôle} La firme miniature affiche des cours et couvre des options sur des plateformes crypto dont les primes sont en bitcoin, et mène des basis trades entre
plateformes crypto et réglementées~; elle a besoin de primes, de deltas et d'une base dans un même jeu de conventions cohérent.
\bfield{Interface} \texttt{coin\_premium(forward, strike, years, vol, right)}, \texttt{usd\_premium}, \texttt{forward\_delta},
\texttt{premium\_adjusted\_delta}~; \texttt{annualised\_basis(future, spot, days, compounding)}~; \texttt{etf\_basis\_carry},
\texttt{return\_on\_margin}~; \texttt{weighted\_median}, \texttt{reference\_rate(trades, start, end, parts)}. Les prix proviennent de
\texttt{firm.parity}.
\bfield{Règles} Les primes en bitcoin sont des prix en dollars divisés par le contrat à terme sous-jacent, jamais par l'indice~; les deltas précisent leur devise~;
la base précise sa capitalisation~; les fenêtres du taux de référence sont semi-ouvertes.
\bfield{Tests d'acceptation} \texttt{code/firm/cryptoopt/tests/}~: une prime en bitcoin face au cours de référence publié par la plateforme~; parité call-put en bitcoin~; deltas à
la monnaie~; annualisation de la base~; arithmétique du carry~; la résistance du taux de référence à une seule transaction importante.
\bfield{Pour aller plus loin} Deltas cohérents avec le sourire de volatilité~; P\&L des contrats à terme margés en crypto, avec marge en bitcoin~; options sur les fonds au comptant avec leurs
conventions en dollars, côte à côte.
\end{build}

\begin{omsources}
\item SEC, «~Statement on the Approval of Spot Bitcoin Exchange-Traded Products~», président G.~Gensler, 10 janvier 2024.
\item iShares Bitcoin Trust ETF, rapport annuel sur Form 10-K pour 2025 (déposé le 27 février 2026), et page web du fonds, septembre 2026.
\item CF Benchmarks, page CME CF Bitcoin Reference Rate, septembre 2026.
\item Deribit, API publique (instruments, synthèses du carnet d'ordres et indice), 24 septembre 2026~; Coinbase Global, rapport annuel sur Form 10-K pour 2025.
\item CFTC, communiqué de presse 7654-17, 1\up{er} décembre 2017~; CME Group, communiqué de presse du 1\up{er} décembre 2017~; SEC, Release No.\ 34-101128, 20 septembre 2024~;
Grayscale Bitcoin Trust, rapport annuel sur Form 10-K pour 2022.
\end{omsources}

\section{Exercices}

\begin{exercise}[$\star$]\label{exo:m3:dated-futures-options-and-etfs:1}
Le cours au comptant vaut 83\,499 et le contrat à terme arrivant à échéance dans 183 jours vaut 85\,669. Calculer la base annualisée, simple et continue.
\end{exercise}

\begin{exercise}[$\star$]\label{exo:m3:dated-futures-options-and-etfs:2}
Un call inverse affiche un cours de référence de 0.0492 bitcoin, le contrat à terme étant à 84\,539. Calculer sa prime en dollars.
\end{exercise}

\begin{exercise}[$\star$]\label{exo:m3:dated-futures-options-and-etfs:3}
Expliquer pourquoi un taux de référence retient des médianes calculées sur des partitions plutôt que la dernière transaction à 4 p.m.
\end{exercise}

\begin{exercise}[$\star\star$]\label{exo:m3:dated-futures-options-and-etfs:4}
Vérifier la parité call-put en bitcoin pour un prix d'exercice de 80\,000, le contrat à terme étant à 84\,539~: le cours de référence du call est de 0.1042 et celui du put de 0.0505.
\end{exercise}

\begin{exercise}[$\star\star$]\label{exo:m3:dated-futures-options-and-etfs:5}
Déterminer le \omterm{def:m3:perpetual-futures:perp}{taux de financement} auquel le trade fonds--contrat à terme atteint son seuil de rentabilité, avec une base de 10~\%, des frais de 0.25~\% et des coûts de 0.30~\% sur les contrats à terme.
\end{exercise}

\begin{exercise}[$\star\star$]\label{exo:m3:dated-futures-options-and-etfs:6}
Expliquer pourquoi la création en nature était plus favorable à l'arbitrage que la \omterm{def:m3:dated-futures-options-and-etfs:cash}{création en espèces}.
\end{exercise}

\begin{exercise}[$\star\star\star$]\label{exo:m3:dated-futures-options-and-etfs:7}
\emph{Programmation.} Avec \texttt{delta\_table}, déterminer les prix d'exercice pour lesquels le delta ajusté de la prime est le plus élevé, et expliquer pourquoi ce n'est pas le
prix d'exercice le plus profond dans la monnaie.
\end{exercise}

\begin{exercise}[$\star\star\star$]\label{exo:m3:dated-futures-options-and-etfs:8}
\emph{Repérer l'erreur.} «~Notre livre de calls inverses est couvert en delta~: nous sommes vendeurs de contrats à terme pour un montant égal à la somme des deltas de Black des calls.~»
\end{exercise}

\section{Problème~: Le basis trade sur ETF}

\begin{problem}[{Problème du week-end --- acheteur du fonds, vendeur du contrat à terme}]\label{pb:m3:dated-futures-options-and-etfs:1}
Un fonds achète pour 200 millions de dollars d'un ETP bitcoin au comptant, financé à 4.5~\%, avec une décote de 10~\%, et vend des contrats à terme réglementés sur le même notionnel
avec une marge de 25~\%. Les frais du fonds s'élèvent à 0.25~\% par an, et les coûts des contrats à terme à 0.30~\% par an. Utiliser le \cref{tab:m3:dated-futures-options-and-etfs:carry}.

\medskip\noindent\textbf{Partie I --- Le carry.}
\begin{enumerate}
\item Avec une base annualisée de 10~\%, que rapporte le trade en un an, en dollars~?
\item Quel capital immobilise-t-il, et quel est son rendement sur capital~?
\item Pour la base instantanée de 4.92~\%, que rapporte-t-il~?
\item Pour quelle base le trade atteint-il son seuil de rentabilité~?
\item Pourquoi la base sur les contrats à terme réglementés est-elle un \omterm{def:m3:perpetual-futures:perp}{taux de financement} déguisé~?
\end{enumerate}

\medskip\noindent\textbf{Partie II --- Le fonctionnement.}
\begin{enumerate}[resume]
\item Comment les parts du fonds restent-elles proches de leur valeur liquidative~?
\item À quel prix le fonds est-il valorisé chaque jour, et comment ce prix est-il calculé~?
\item À quel prix le contrat à terme est-il réglé~?
\item Pourquoi les deux peuvent-ils différer à l'échéance, et que cela ferait-il au trade~?
\item Qu'est-ce qui a changé le 29 juillet 2025, et pourquoi cela compte-t-il pour ce trade~?
\end{enumerate}

\medskip\noindent\textbf{Partie III --- Les risques.}
\begin{enumerate}[resume]
\item La base s'élargit de 10~\% à 20~\% un mois après l'entrée. Que deviennent la position et sa marge~?
\item Le financeur relève la décote à 25~\%. Que devient le rendement sur capital~?
\item Le fonds se traite avec une décote de 1~\% le jour où le trade est dénoué. Que coûte cela~?
\item Que devient le trade si le bitcoin chute de 30~\%~?
\item Comment ce même trade se présenterait-il avec un contrat à terme sur crypto margé en crypto plutôt qu'avec un contrat réglementé~?
\end{enumerate}

\medskip\noindent\textbf{Partie IV --- Jugement.}
\begin{enumerate}[resume]
\item Qui se trouve de l'autre côté du trade, en payant la base~?
\item Pourquoi un grand nombre de fonds menant ce trade en même temps peuvent-ils rendre le marché fragile~?
\item Ce trade est-il un arbitrage~?
\item Énoncer le \emph{résultat nommé}~: le carry annualisé du trade et son rendement sur capital pour une base de 10~\%, ainsi que son carry pour la base instantanée
~.
\item En une phrase~: qu'ont apporté les fonds au comptant au marché de la crypto~?
\end{enumerate}
\end{problem}

\section{Questions d'entretien}

\begin{interviewq}[$\star$ \iqroles{trader}]\label{iq:m3:dated-futures-options-and-etfs:1}
Qu'est-ce qu'une \omterm{def:m3:dated-futures-options-and-etfs:inverse}{option inverse}, et comment convertissez-vous sa prime en dollars~?
\end{interviewq}

\begin{interviewq}[$\star$ \iqroles{chercheur}]\label{iq:m3:dated-futures-options-and-etfs:2}
Comment un taux de référence bitcoin est-il construit, et pourquoi de cette manière~?
\end{interviewq}

\begin{interviewq}[$\star\star$ \iqroles{trader}]\label{iq:m3:dated-futures-options-and-etfs:3}
Pourquoi le delta ajusté de la prime d'un call inverse profond dans la monnaie diminue-t-il à mesure qu'on s'enfonce davantage~?
\end{interviewq}

\begin{interviewq}[$\star\star$ \iqroles{trader, risque}]\label{iq:m3:dated-futures-options-and-etfs:4}
Expliquez le basis trade entre l'ETF au comptant et les contrats à terme réglementés, ainsi que ses risques.
\end{interviewq}

\begin{interviewq}[$\star\star$ \iqroles{chercheur}]\label{iq:m3:dated-futures-options-and-etfs:5}
En quoi la surface de volatilité du bitcoin diffère-t-elle d'une surface sur indice actions~?
\end{interviewq}

\begin{interviewq}[$\star\star\star$ \iqroles{développeur}]\label{iq:m3:dated-futures-options-and-etfs:6}
Concevez un système de gestion des risques pour un livre d'options comprenant des \omterm{def:m3:dated-futures-options-and-etfs:inverse}{options inverses} sur une plateforme et des options en dollars sur une autre.
\end{interviewq}
